% arXiv submission — Efficient Mechanistic Circuit Analysis of Open-Weight LLMs on Apple Silicon
% J. Melton — July 2026
%
% OPEN ITEMS BEFORE FINAL SUBMISSION:
%   1. Resolve L24H15 vs. L24H16 discrepancy (cross-model-comparison.md vs. results-summary.md ablation rank table).
%   2. Standardize L1H11 CI upper bound (+0.005 or +0.006); verify against data/ioi/statistical-validation.json.
%   4. Table 2 is ordered by ablation drop, not combined rank — caption notes this but consider reordering.
%   5. Confirm SwiftSci Interp repository URL and insert below (search for REPO_URL_HERE).
%   6. Collect and add efficiency benchmark data for Section 6.
%   7. Complete path-patching experiments (Section 3.4) and revise Section 6/7 accordingly.

\documentclass[11pt,a4paper]{article}

% ── Encoding & fonts ──────────────────────────────────────────────────────────
\usepackage[utf8]{inputenc}
\usepackage[T1]{fontenc}
\usepackage{lmodern}

% ── Math ──────────────────────────────────────────────────────────────────────
\usepackage{amsmath}
\usepackage{amssymb}

% ── Layout ────────────────────────────────────────────────────────────────────
\usepackage[margin=1.25in]{geometry}
\usepackage{microtype}
\usepackage{parskip}
\usepackage{setspace}

% ── Figures & tables ──────────────────────────────────────────────────────────
\usepackage{graphicx}
\usepackage{booktabs}
\usepackage{float}
\usepackage{caption}
\usepackage{subcaption}
\usepackage{array}
\usepackage{multirow}

% ── Code listings ─────────────────────────────────────────────────────────────
\usepackage{listings}
\usepackage{xcolor}

\definecolor{codegray}{rgb}{0.97,0.97,0.97}
\definecolor{commentgreen}{rgb}{0.0,0.5,0.0}
\definecolor{keywordblue}{rgb}{0.0,0.0,0.7}
\definecolor{stringbrown}{rgb}{0.5,0.1,0.0}

\lstdefinestyle{python}{
  language=Python,
  backgroundcolor=\color{codegray},
  basicstyle=\small\ttfamily,
  keywordstyle=\color{keywordblue}\bfseries,
  commentstyle=\color{commentgreen},
  stringstyle=\color{stringbrown},
  numberstyle=\tiny\color{gray},
  numbers=left,
  numbersep=6pt,
  stepnumber=1,
  frame=single,
  framerule=0.4pt,
  xleftmargin=2em,
  breaklines=true,
  showstringspaces=false,
  tabsize=4,
  captionpos=b,
}

% ── Citations ─────────────────────────────────────────────────────────────────
\usepackage[numbers,sort&compress]{natbib}

% ── Hyperlinks ────────────────────────────────────────────────────────────────
\usepackage[colorlinks=true,linkcolor=blue,citecolor=blue,urlcolor=blue,
            pdftitle={Efficient Mechanistic Circuit Analysis of Open-Weight LLMs on Apple Silicon},
            pdfauthor={J. Melton}]{hyperref}

% ─────────────────────────────────────────────────────────────────────────────

\title{\textbf{Efficient Mechanistic Circuit Analysis of\\Open-Weight LLMs on Apple Silicon}}

\author{
  J.\ Melton \\[4pt]
  Independent Researcher \\
  \texttt{jacob@paydirt.solutions}
}

\date{July 2026 \\ \smallskip \small Preprint. Under review.}

% ─────────────────────────────────────────────────────────────────────────────
\begin{document}
\maketitle

% ── Abstract ──────────────────────────────────────────────────────────────────
\begin{abstract}
Mechanistic interpretability research has largely depended on GPU clusters, limiting
reproducibility and iteration speed for many researchers. We demonstrate that full
activation-patching and mean-ablation circuit analyses of multi-billion-parameter
transformer models are tractable on Apple Silicon using the MLX framework, with no
specialized hardware beyond a modern Mac. Applying this infrastructure to the
Indirect Object Identification (IOI) task on Llama-3.2-3B and Pythia-1.4B, we find a
pronounced bottleneck structure in both models: a single dominant head (L15H20 in
Llama, L10H7 in Pythia) accounts for 22--28\% of clean logit difference on its own,
substantially more concentrated than the distributed 26-component circuit Wang et al.\
(2022) report in GPT-2 Small. Both circuits occupy three depth clusters broadly
matching Wang et al.'s functional zones.

We also report a negative methodological result. Matching circuit-critical head
positions between the two models by relative depth yields 8 of 10 pairs within
$\pm$0.075---but under a permutation null, randomly chosen head sets match at 59\%
(uniform null, $p = 0.098$) or 69\% (null restricted to the layer band where critical
heads actually sit, $p = 0.32$). The 80\% figure is not distinguishable from chance,
and the matching is functionally uninformative besides: it pairs Llama's dominant head
with Pythia's weakest. We report the null because the depth-conservation claim is the
one a reader is most likely to take away, and it does not hold.

Path patching over the five highest-scoring Llama heads shows that L15H20 receives a
substantial direct contribution from L14H0 and is largely independent of the later
heads, consistent with its role as an early bottleneck. Factual-association patching
recovers heads in five of the same layers as IOI but at different head indices. Our
positive claims are that IOI admits a sparse, bottlenecked circuit in both models, and
that Apple Silicon is a viable substrate for reproducible mech-interp research at the
1--3B parameter scale.
\end{abstract}

% ─────────────────────────────────────────────────────────────────────────────
\section{Introduction}
\label{sec:intro}

Mechanistic interpretability aims to reverse-engineer the algorithms implemented by
neural networks---to move from ``the model predicts X'' to ``these specific
components, performing these specific computations, cause the model to predict X.''
The dominant paradigm for this reverse-engineering is \emph{circuit analysis}:
identify the smallest subset of attention heads and MLP layers that is both causally
necessary and causally sufficient to reproduce a target behavior, then assign
functional roles to each component within that subset.

The foundational work in this paradigm---Wang et al.\ (2022)---demonstrated that the
Indirect Object Identification (IOI) task in GPT-2 Small is implemented by a
comprehensible circuit of 26 attention heads organized into three functional classes:
duplicate-token heads that flag repeated names, S-inhibition heads that suppress the
subject position, and name-mover heads that copy the indirect object name to the
output. The result was striking: a clean, human-interpretable algorithm, found inside
a neural network trained only to predict next tokens.

What remains unresolved is whether this algorithm is specific to GPT-2 Small, or
whether it reflects something deeper---a functional structure that emerges wherever a
language model learns to solve this problem. GPT-2 Small is small (117M parameters),
trained on a narrow corpus (WebText), and architecturally plain (standard MHA, learned
positional embeddings). Modern open-weight models differ on every dimension: they are
larger by one to two orders of magnitude, trained on multitrillion-token corpora, and
equipped with grouped-query attention, rotary position embeddings, and gated MLP
activations. If the IOI circuit is merely an artifact of GPT-2's architecture, we
would expect a different circuit in a different model. If it reflects a general
algorithmic solution, we would expect the same functional structure---with head indices
shifted, but depth relationships and role assignments preserved.

This paper answers that question by replicating and extending Wang et al.'s analysis
on two modern models: Llama-3.2-3B \citep{meta2024llama3} (a modern base model with
grouped-query attention and rotary embeddings; we use the base weights,
\texttt{mlx-community/Llama-3.2-3B-bf16}, not the instruction-tuned variant) and Pythia-1.4B
\citep{biderman2023pythia} (a fully documented, standard-MHA model trained on The
Pile, making it the closest modern analog to GPT-2 Small's controlled training
conditions). Using activation patching and mean ablation, we identify circuit-critical
heads in each model, compare their positions to Wang et al.'s findings, and test for
cross-architecture conservation at the level of relative network depth rather than
absolute head indices.

\paragraph{Apple Silicon as interpretability infrastructure.}
All experiments in this paper were run on Apple Silicon hardware using the MLX
framework, without access to NVIDIA GPU clusters. This is a deliberate choice. The
interpretability tooling ecosystem---TransformerLens in particular---assumes CUDA, and
the community's empirical norms have drifted toward hardware accessible only to
well-resourced labs. Apple's M-series chips offer up to 96\,GB of unified memory
shared between CPU and GPU, which is well-suited to the memory access patterns of
activation patching: you must store one full forward pass of cached activations per
example, then perform hundreds of re-runs patching each (layer, head, token position)
triple. On an M-series machine with sufficient unified memory, a 3B-parameter model
fits in full float16 with room for the activation cache, and forward passes are fast
enough that the full IOI patching sweep over 100 examples completes in wall-clock hours
rather than days. We describe our Apple Silicon-compatible pipeline and release it
alongside this paper to lower the barrier for researchers without cluster access to run
mechanistic interpretability experiments at the 1--7B parameter scale.

\subsection*{Contributions}

\begin{enumerate}
  \item \textbf{Cross-architecture circuit replication.}  We identify circuit-critical
    heads for the IOI task in Llama-3.2-3B and Pythia-1.4B using activation patching
    (sufficiency) and mean ablation (necessity). In both models, a small bottleneck
    structure emerges: a single dominant head accounts for 22--28\% of the model's
    logit difference for the IO token, substantially more concentrated than the
    distributed 26-component circuit in GPT-2 Small.

  \item \textbf{A permutation null for depth-based cross-architecture matching.}
    Circuit-critical heads in both models cluster into three depth bands broadly
    compatible with Wang et al.'s functional zones. But the specific quantitative
    claim that invites---that 80\% of head positions are conserved within $\pm$0.075
    relative depth---does not survive a null test: random head sets match at
    59--69\% under the same greedy procedure ($p = 0.098$ and $p = 0.32$ under two
    nulls). We report the null, the reason the tolerance is too permissive to be
    informative, and what a valid test of depth conservation would require.

  \item \textbf{Edge-level path patching of the Llama bottleneck.}  We path-patch the
    ten sender--receiver edges among Llama's five highest-scoring heads, finding that
    L15H20's dominance is fed substantially by L14H0 (normalized score 0.165) and is
    largely decoupled from the later heads.

  \item \textbf{Layer-level overlap across tasks.}  Factual-association patching in
    Llama-3.2-3B recovers top-5 heads occupying five of the same layers as five of the
    top-9 IOI heads, at different head indices. We report the associated null: with
    both tasks' heads concentrated in the same mid-to-late band, this degree of layer
    overlap is not by itself strong evidence of shared topology.

  \item \textbf{An Apple Silicon-compatible activation patching pipeline.}  We release
    tooling built on the MLX framework that replicates TransformerLens's hook-based
    patching interface for LlamaForCausalLM and GPTNeoXForCausalLM models, enabling
    circuit analysis on Apple Silicon hardware without modification.
\end{enumerate}

% ─────────────────────────────────────────────────────────────────────────────
\section{Related Work}
\label{sec:related}

\subsection{A Mathematical Framework for Transformer Circuits}

The theoretical foundations of circuit analysis in transformers were laid by Elhage et
al.\ (2021) \citep{elhage2021mathematical}, who introduced the \emph{transformer
circuits} framework. The key insight is that a transformer's residual stream can be
analyzed as a sum of contributions from attention heads and MLP layers, with each
attention head performing a rank-one update of the residual: the head reads from the
stream via query and key projections, computes a weighted average of value projections,
and writes back a low-rank update whose structure is determined by the composition of
the OV and QK matrices.

Elhage et al.\ used this decomposition to characterize \emph{induction heads}---a
two-head circuit found across models that implements the copying operation ``if token A
was followed by token B earlier in context, predict B again when A reappears.''
Induction heads are relevant to IOI analysis because the same copying mechanism
underlies name-mover head behavior: identifying the indirect object earlier in context
and copying its token to the output position. Our results confirm that the depth zone
corresponding to induction-like behavior (relative depth 0.40--0.55) is present in
both Llama-3.2-3B and Pythia-1.4B.

The circuits framework also establishes \emph{virtual attention heads}---paths in
which the output of one head is read as part of another head's input---as the primary
mechanism for inter-head communication. This framing is a prerequisite for
understanding the IOI circuit's S-inhibition subgraph, where a set of early heads
writes information about the subject's position that is then used by a later head to
suppress subject-token logits.

\subsection{Interpretability in the Wild: IOI in GPT-2 Small}

Wang et al.\ (2022) \citep{wang2022interpretability} is the direct predecessor to this
work. Using the IOI task---sentences of the form ``When Mary and John went to the
store, John gave a drink to \underline{\phantom{Mary}}''---they applied activation
patching across all attention heads and MLP layers in GPT-2 Small (12 layers, 12
heads) to identify which components causally mediate correct IO prediction. The result
was a 26-component circuit with three functional classes:

\begin{itemize}
  \item \textbf{Name-mover heads} (L9H6, L9H9, L10H0; relative depths 0.75--0.83)
    copy the indirect object name to the logit output at the final token position.
  \item \textbf{S-inhibition heads} (L7H3, L7H9, L8H6, L8H10; relative depths
    0.58--0.67) suppress the subject token's representation in the residual stream,
    preventing the model from predicting the repeated name.
  \item \textbf{Duplicate-token heads and induction heads} (L0H1, L3H0, L5H5, L5H8;
    relative depths 0.0--0.42) flag the position of the repeated subject token so
    that downstream heads know which name to suppress.
\end{itemize}

Our work replicates this three-zone structure across Llama-3.2-3B and Pythia-1.4B,
extending the generalizability claim beyond GPT-2 Small. We also find that in larger
models the circuit has a more concentrated ``bottleneck'' structure, with one head
contributing disproportionately to logit difference---a difference from GPT-2 Small's
more distributed allocation that we discuss in Section~\ref{sec:comparison}.

\subsection{Circuit Analysis Beyond GPT-2 Small}
\label{sec:related-scale}

Several lines of work bear directly on the questions this paper asks, and situate our
contribution more narrowly than an unqualified ``does IOI generalize?'' framing would
suggest.

\citet{lieberum2023scale} conduct circuit analysis on Chinchilla-70B, finding that
activation patching and attention-pattern analysis remain tractable at that scale and
that identifiable circuits exist for multiple-choice tasks. Their result establishes
that circuit analysis scales; ours addresses the complementary question of whether a
\emph{specific} circuit recurs across architectures.

\citet{merullo2024reuse} is the closest prior work to our cross-task experiment
(\S\ref{sec:cross-task}). They show that the IOI circuit in GPT-2 Medium is partially
reused for a Colored Objects task, with overlapping heads playing analogous roles.
Their evidence for cross-task reuse is stronger than ours: they establish role
correspondence via path patching and targeted interventions, whereas we report layer
co-occurrence. Readers interested in cross-task circuit reuse should treat their
result as primary and ours as a weak replication in a different model.

\citet{conmy2023acdc} introduce ACDC, which automates circuit discovery by iteratively
pruning edges from the full computational graph against a faithfulness threshold. ACDC
is the natural baseline for the head-selection procedure in \S\ref{sec:stats}, and we
do not compare against it---our selection is a hand-specified ranking over marginal
effects, which is simpler but has no faithfulness guarantee. \citet{hanna2023greaterthan}
apply similar methods to arithmetic, further mapping the space of known circuits.

On methodology, \citet{zhang2024bestpractices} and \citet{heimersheim2024patching}
document how sensitive activation patching results are to choices we make without much
justification: the corruption distribution, the metric (logit difference versus
probability versus KL), and whether patching is done at the head output or at
individual Q/K/V projections. Our protocol---name-swap corruption, normalized logit
difference, head-output patching---is a common configuration but not the only defensible
one, and we did not test robustness to these choices.

\subsection{TransformerLens}

TransformerLens \citep{nanda2022transformerlens} is the standard tooling for
mechanistic interpretability experiments on transformer language models. It provides a
hook-based interface for intercepting and modifying activations at any point in the
forward pass---attention head outputs, MLP outputs, residual stream positions, or
individual Q/K/V projections---making activation patching experiments straightforward
to implement.

Our work builds on TransformerLens for Pythia-1.4B, where the library's support is
mature. For Llama-3.2-3B with grouped-query attention, we implement a compatible
interface that maps the 24-query-head, 8-KV-head structure to a consistent
per-query-head naming convention. GQA complicates patching because each KV head is
shared across a group of 3 query heads; we patch at the query-head output level, after
the attention weights are applied to the shared value projections, to preserve
head-level attribution granularity. We release this extension alongside the paper.

\subsection{Locating and Editing Factual Associations}

Meng et al.\ (2022) \citep{meng2022locating} introduced ROME (Rank-One Model Editing),
which demonstrated that factual associations---the model's belief that ``the Eiffel
Tower is in Paris''---are stored with high locality in mid-layer MLP weights, and can
be surgically edited by a rank-one update to those weights. The key interpretability
contribution was not the editing procedure itself but the preceding localization
experiment: using causal tracing (a form of activation patching), they showed that
early-layer attention heads read subject tokens and route information to mid-layer
MLPs, which then store and retrieve the associated object.

This causal tracing methodology is closely related to our activation patching protocol.
A key difference is that ROME's target is stored parametric knowledge (factual
associations mediated by MLP weights), while IOI involves in-context reasoning (copying
from the prompt, not from stored facts). Our factual association experiments in
Section~\ref{sec:cross-task} partially bridge this gap: we apply patching to factual
association prompts and compare the resulting head-importance rankings to the IOI
results, finding that layer-level circuit structure is shared while head-level
assignments diverge. This cross-task comparison was not reported in ROME and provides
new evidence about how circuits for different types of factual reasoning coexist within
the same model.

\subsection{Towards Monosemanticity and Sparse Autoencoders}

Elhage et al.\ (2022) \citep{elhage2022toy} and the Anthropic mechanistic
interpretability team have pursued an orthogonal but complementary program: rather than
identifying circuits for specific behaviors, they seek to decompose the residual stream
into \emph{features}---directions in activation space that correspond to
human-interpretable concepts---using sparse autoencoders (SAEs). The central finding
of the monosemanticity work is that individual neurons are typically polysemantic
(responding to unrelated concepts), but a sparse linear decomposition of neuron
activations recovers interpretable monosemantic features.

The circuit-tracing and monosemanticity programs address different levels of the
interpretability problem. Circuit analysis asks: which components mediate this
behavior? Feature decomposition asks: what information is represented in the
activations those components read and write? They are complementary: knowing that
L15H20 in Llama-3.2-3B is the dominant IOI head is more interpretable if we also know
what feature L15H20 reads from the residual stream (presumably a representation of the
indirect object token's identity and position). Our work focuses on the circuit level,
but Section~\ref{sec:discussion} discusses how SAE-based feature attribution could be
applied to the circuit-critical heads we identify, and we propose this as a natural
extension.

% ─────────────────────────────────────────────────────────────────────────────
\section{Methods}
\label{sec:methods}

\subsection{Experimental Infrastructure: SwiftSci Interp}

All experiments were conducted using SwiftSci Interp, an activation-patching harness
for Apple Silicon built on the MLX framework.\footnote{SwiftSci Interp is available
at \url{https://github.com/jmelton-research/swiftsci-interp}.
} The library implements a hook-compatible interface analogous to TransformerLens's
activation interception API, targeting MLX-native model implementations loaded via
\texttt{mlx-lm}. Its core design principle is that every attention module in a model
is hot-swapped at load time with a patchable drop-in replacement that shares the
original weight tensors---no copies---and adds a mode-switched interception point
around the scaled dot-product attention (SDPA) computation.

\paragraph{Patchable module design.}
For Llama-3.2-3B, the drop-in is \texttt{PatchableAttention}, which wraps the
original \texttt{LlamaAttention} module. For Pythia-1.4B,
\texttt{PythiaPatchableAttention} wraps the \texttt{GPTNeoXAttention} module, routing
through the fused \texttt{query\_key\_value} projection to reconstruct per-head Q, K,
V tensors before the interception point. Both modules support three operating modes:

\begin{itemize}
  \item \texttt{normal}---identical behavior to the original module; introduces no
    overhead beyond the mode check.
  \item \texttt{cache\_clean} / \texttt{cache\_corrupt}---stores the full SDPA output
    tensor (shape $[B, n_\text{heads}, L, d_\text{head}]$) to a per-module buffer
    during the forward pass. Used to populate the clean and corrupt activation caches
    without a separate caching pass.
  \item \texttt{patch}---runs the forward pass on the input normally, then replaces
    one head slice in the SDPA output with the corresponding slice from the clean cache
    before passing through the output projection.
\end{itemize}

\paragraph{GQA handling.}
Llama-3.2-3B uses grouped-query attention (GQA) \citep{ainslie2023gqa} with 24 query
heads and 8 KV heads. The patching interception occurs after the attention weights are
applied to the (broadcast) value projections, at the query-head output level. This
gives a 24-element patch space that mirrors the query-head count, preserving head-level
attribution granularity. Patching at the KV-head level would conflate the
contributions of the three query heads that share each KV pair; patching at the
query-head SDPA output avoids this confound. We verified that in \texttt{normal} mode
both patchable modules reproduce the original model's logits on a held-out test input
to within floating-point precision before running any experimental passes.

\paragraph{GQA key-value broadcast.}
During the corrupt or clean pass, the values broadcast from 8 KV heads to 24 query
heads occur inside \texttt{scaled\_dot\_product\_attention}; the SDPA buffer therefore
stores the full 24-head output, not the 8-head compressed form. This means each of the
24 patch positions is truly independent, and patching query-head $h$ does not
implicitly alter the other two heads sharing the same KV pair.

\paragraph{Pythia module layout.}
MLX-LM exposes Pythia's layers as \texttt{model.layers[i].attention} rather than
\texttt{model.model.layers[i].self\_attn}. The patchable replacement reuses the
\texttt{dense} (output projection) attribute rather than \texttt{o\_proj}. All other
behavior is identical; the mode-switch logic is shared.

\subsection{Activation Patching Protocol}
\label{sec:patching}

Activation patching (also called causal tracing when applied to individual stream
positions) answers the question of \emph{sufficiency}: if a component's activations in
the clean run are inserted into the corresponding position of a corrupted run, does the
model's output recover toward the clean prediction? A high recovery score indicates
that the component is sufficient to carry the relevant information on its own.

\paragraph{IOI corruption scheme.}
For each example $(S, \mathit{IO}, \text{prompt})$, the corrupted version swaps the
subject and indirect object names: the corrupted prompt is constructed from the same
template with $S_\text{corrupt} = \mathit{IO}$ and $\mathit{IO}_\text{corrupt} = S$.
Because the two name slots have the same token count for all examples in our dataset
(a requirement enforced at dataset generation time), clean and corrupt sequences are
token-position aligned. This ensures that cached SDPA tensors from the clean run can
be substituted at the same token positions in the corrupt run without shape mismatch.

\paragraph{Metric.}
The target metric is logit difference (LD):
$\text{LD} = \text{logit}(\mathit{IO}) - \text{logit}(S)$ at the final token position.
Positive LD indicates the model favors the IO name over the subject name. The
\emph{normalized logit-difference recovery} for head $(l, h)$ on example $i$ is:

\begin{equation}
  r_{l,h}^{(i)} = \frac{
    \text{LD}^{(i)}_{\text{patch}(l,h)} - \text{LD}^{(i)}_{\text{corrupt}}
  }{
    \text{LD}^{(i)}_{\text{clean}} - \text{LD}^{(i)}_{\text{corrupt}}
  }
  \label{eq:recovery}
\end{equation}

where $\text{LD}^{(i)}_{\text{patch}(l,h)}$ is the logit difference of a forward pass
on the corrupt input in which only head $h$ at layer $l$ has been replaced by its
clean-run SDPA output. A score of 1.0 means full recovery to the clean logit
difference; 0.0 means the patch had no effect; values below 0 or above 1 indicate
interference effects. The reported score for each head is
$\hat{r}_{l,h} = \frac{1}{n}\sum_{i=1}^n r_{l,h}^{(i)}$, the mean over all $n$
examples.

\paragraph{Pass structure.}
A full patching sweep requires $2 + n_L \times n_H$ forward passes over the example
batch:

\begin{enumerate}
  \item \textbf{Clean pass} --- all modules in \texttt{cache\_clean} mode; SDPA
    outputs are stored and logit differences are computed.
  \item \textbf{Corrupt pass} --- all modules in \texttt{cache\_corrupt} mode; SDPA
    outputs are stored and baseline corrupt logit differences are computed.
  \item \textbf{Sweep} --- for each $(l, h)$: all modules in \texttt{normal} mode;
    module $l$ is set to \texttt{patch} mode with \texttt{patch\_head = h}; one
    forward pass on the corrupt input recovers
    $\text{LD}^{(i)}_{\text{patch}(l,h)}$ for all $i$ simultaneously.
\end{enumerate}

For Llama-3.2-3B this yields $2 + 28 \times 24 = 674$ total forward passes; for
Pythia-1.4B, $2 + 24 \times 16 = 386$ passes. All passes are fully batched over all
$n = 100$ examples.

\begin{lstlisting}[style=python, caption={Pseudocode for the activation patching sweep. \texttt{attns} is a list of \texttt{PatchableAttention} modules, one per layer. \texttt{set\_all\_modes} sets every module's mode simultaneously.}]
# Inputs:
#   attns: list of PatchableAttention, one per layer (len = n_layers)
#   clean_ids:   [n_examples, seq_len]  -- clean tokenized prompts
#   corrupt_ids: [n_examples, seq_len]  -- IO/S-swapped prompts

# Pass 1: cache clean activations
set_all_modes(attns, mode="cache_clean")
clean_logits = model(clean_ids)
eval(clean_logits, *(a.clean_sdpa for a in attns))
clean_ld = logit_diff(clean_logits, io_ids, s_ids)  # [n_examples]

# Pass 2: cache corrupt activations
set_all_modes(attns, mode="cache_corrupt")
corrupt_logits = model(corrupt_ids)
eval(corrupt_logits, *(a.corrupt_sdpa for a in attns))
corrupt_ld = logit_diff(corrupt_logits, io_ids, s_ids)

# Passes 3..674: patching sweep
scores = zeros([n_layers, n_heads])
for l in range(n_layers):
    for h in range(n_heads):
        set_all_modes(attns, mode="normal")
        attns[l].mode = "patch"
        attns[l].patch_head = h       # only layer l patches head h
        patched_logits = model(corrupt_ids)
        eval(patched_logits)
        patched_ld = logit_diff(patched_logits, io_ids, s_ids)
        denom = clean_ld - corrupt_ld               # [n_examples]
        recovery = (patched_ld - corrupt_ld) / denom
        scores[l, h] = mean(recovery)               # scalar
\end{lstlisting}

\textbf{GQA note.} \texttt{attns[l].patch\_head = h} indexes into the 24-dimensional
query-head axis regardless of the 8-KV-head layout. The \texttt{clean\_sdpa} buffer
holds the full $[B, 24, L, d_\text{head}]$ output; the slice
\texttt{clean\_sdpa[:, h:h+1, :, :]} is substituted.

\subsection{Mean Ablation Protocol}
\label{sec:ablation}

Mean ablation addresses \emph{necessity}: does removing a component's
contribution---replacing it with a neutral baseline---degrade the model's performance
on the task? Where activation patching asks whether a component \emph{can} carry the
relevant information, mean ablation asks whether the model \emph{relies} on it in
normal operation.

\paragraph{Reference distribution.}
The ablation target for head $(l, h)$ is the mean SDPA output of that head computed
over all $n = 100$ clean IOI examples:
$\bar{A}_{l,h} = \frac{1}{n}\sum_{i=1}^n A_{l,h}^{(i)}$, where
$A_{l,h}^{(i)} \in \mathbb{R}^{L \times d_{\text{head}}}$ is the per-example SDPA
output at each token position. The mean is taken over the batch dimension, preserving
the sequence-position structure. Ablating with the mean over the same distribution on
which the model is evaluated (rather than, for example, a random Gaussian or a
separately collected reference set) ensures that the ablated output is a plausible
activational state for the model rather than an out-of-distribution injection.

\paragraph{Metric.}
The \emph{ablation drop} for head $(l, h)$ is:

\begin{equation}
  \Delta_{l,h} = \overline{\text{LD}}_{\text{clean}} - \overline{\text{LD}}_{\text{ablate}(l,h)}
  \label{eq:ablation}
\end{equation}

where $\overline{\text{LD}}_{\text{clean}}$ is the mean clean logit difference over
all examples and $\overline{\text{LD}}_{\text{ablate}(l,h)}$ is the mean logit
difference when head $(l, h)$'s SDPA output is replaced by $\bar{A}_{l,h}$ in every
example. Positive $\Delta_{l,h}$ indicates the head boosts the IO logit relative to
the subject; negative values indicate suppression of correct IO prediction (which can
reflect an interference-suppression role rather than direct IO promotion).

\paragraph{Necessity threshold.}
A head is designated circuit-critical under the necessity criterion when its ablation
drop exceeds 20\% of the clean mean logit difference. For Llama-3.2-3B
($\overline{\text{LD}}_{\text{clean}} = 5.643$) this corresponds to
$\Delta_{l,h} > 1.130$; for Pythia-1.4B ($\overline{\text{LD}}_{\text{clean}} =
4.120$) the threshold is $\Delta_{l,h} > 0.824$. In practice the necessity threshold
is used as an interpretive anchor rather than a hard inclusion criterion; circuit
membership is determined by the combined rank score described in Section
\ref{sec:stats}.

\paragraph{Pass structure.}
The ablation sweep requires $1 + n_L \times n_H$ forward passes: one clean pass to
cache per-head means, then one ablation pass per $(l, h)$ pair. For Llama this is
$1 + 672 = 673$ total passes, comparable to the patching sweep.

\subsection{Path Patching for Causal Verification}
\label{sec:path-patching}

Activation patching and mean ablation identify individual heads as causally relevant,
but they do not distinguish \emph{direct} contributions from \emph{mediated}
contributions. A head may score highly on both metrics because it directly writes to
the logit output at the final position, or because it is an intermediary in a longer
causal path---for example, writing to the residual stream at an intermediate token
position that is later read by a downstream head that writes to the final position.

Path patching \citep{goldowsky2023localizing} extends activation patching to directed
paths between specific components. Rather than patching a single component's output, a
path patch replaces the clean-run activation that flows along one specific edge in the
computational graph---from the output of component $A$ as read by the input of
component $B$, while holding all other paths fixed. A path that shows high normalized
recovery under path patching is causally necessary specifically through the direct
$A \to B$ connection, ruling out the mediated-path alternative.

In our setting, path patching is used to characterize how the dominant head identified
by activation patching (L15H20 in Llama) relates to the other high-scoring heads. Path
patching is implemented within SwiftSci Interp by extending the \texttt{patch} mode to
accept a target-receiver argument.

Each edge is measured with a four-pass protocol: a clean baseline; a corrupted run
capturing the sender's activation; an intermediate clean run with the sender patched,
capturing the receiver's resulting input; and a final clean run with the receiver's
input replaced. The score is the normalized logit-difference change,
$(\mathrm{LD}_{\text{clean}} - \mathrm{LD}_{\text{patched}}) /
(\mathrm{LD}_{\text{clean}} - \mathrm{LD}_{\text{corrupt}})$, over the same
$n = 100$ IOI examples, with
$\mathrm{LD}_{\text{clean}} = 5.668$ and
$\mathrm{LD}_{\text{corrupt}} = -5.754$. We measured all ten ordered sender--receiver
edges among Llama's five highest-scoring heads (L14H0, L15H20, L17H17, L19H1, L24H15);
results are reported in \S\ref{sec:path-results}. Pythia was not path-patched, and the
role of L10H7 therefore remains uncharacterized at the edge level.

\subsection{Statistical Validation}
\label{sec:stats}

\paragraph{Bootstrap confidence intervals.}
For all heads meeting the inclusion criterion (top-10 by combined rank score, or mean
patching score $\geq 0.03$), we compute 95\% bootstrap confidence intervals on the
mean patching score. The bootstrap resamples the $n = 100$ per-example recovery scores
with replacement, generating 1{,}000 resampled datasets and computing the mean for
each. The 95\% CI is the $[2.5^{\text{th}}, 97.5^{\text{th}}]$ percentile of the
bootstrap distribution of the mean. Resampling is seeded with seed 42 for
reproducibility. Full CI data: \texttt{data/ioi/statistical-validation.json}.

\paragraph{Significance criterion.}
A head's patching score is considered statistically reliable if its 95\% CI excludes
zero. This corresponds to the claim that the head's mean recovery score is positive in
expectation across the example distribution, not merely a noise artifact of the
particular 100 examples observed. All top-10 heads in both models satisfy this
criterion with substantial margin: for the top three heads in each model, the CI lower
bound exceeds the CI width by a factor of at least 3.

\paragraph{Anomalous patching scores.}
Pythia-1.4B's L1H11 presents a dissociation: its mean patching score is slightly
negative ($-0.003$, 95\% CI $[-0.011, +0.005]$, straddling zero), yet its ablation
drop is substantial ($0.327$, rank 4 by ablation). This dissociation is consistent
with a suppression role: the head reduces logit difference when ablated (its absence
degrades performance), but restoring its clean activation in a corrupt context does not
recover logit difference (it cannot by itself make the model prefer IO over S). This
pattern is characteristic of duplicate-token or induction heads that gate downstream
components rather than writing directly to the output logits. L1H11 is excluded from
all patching-based analyses and retained in cross-model analysis solely on the basis of
its ablation score.

\paragraph{Cross-model comparison: combined rank score.}
To compare head importance across models that differ in scale, architecture, and
absolute logit-difference baselines, we compute a combined normalized rank score. For
each model separately, both the patching scores and the ablation drop scores are
min-max normalized to $[0, 1]$ across all $(l, h)$ pairs. The combined score is the
sum of the two normalized ranks. The top-10 heads by combined score are reported as
circuit-critical for each model. Ties within the combined score are broken by patching
rank.

\paragraph{Cross-architecture position matching.}
For the cross-model comparison, each head is mapped to a scalar \emph{relative depth}
$d = l / n_L$, where $l$ is the zero-indexed layer number and $n_L$ is the total
number of layers. Positions are matched greedily in ascending order of $|\Delta d|$
between the two models' top-10 head sets; a match is accepted if $|\Delta d| \leq
0.075$. This tolerance was chosen to be smaller than the spacing between the four
functional depth zones identified by Wang et al.\ (2022) (${\approx}0.15$--$0.20$
between zone midpoints), ensuring that zone-level matches are not spuriously collapsed
across zone boundaries.

\subsection{Dataset}
\label{sec:dataset}

The IOI dataset consists of $n = 100$ examples drawn from the template:

\begin{quote}
  ``When \{S\} and \{IO\} went to the store, \{S\} gave a bottle to''
\end{quote}

where $S$ (subject / repeated name) and $\mathit{IO}$ (indirect object / target name)
are sampled from a pool of common English given names, matched to ensure that both
names tokenize to exactly one token under each model's tokenizer. All 100 examples
satisfy the single-token constraint for both Llama-3.2-3B and Pythia-1.4B; examples
failing this constraint were rejected at generation time. Corrupted prompts are
produced by swapping $S$ and $\mathit{IO}$ within the same template, preserving
sequence length. Full dataset: \texttt{data/ioi/dataset.json}.

The factual association dataset consists of $n = 50$ hand-constructed
subject--relation--object triples across two relation templates: 25 \texttt{capital\_of}
items (``The capital of \{subject\} is'') and 25 \texttt{official\_language\_of} items
(``The official language of \{subject\} is''). Subjects are countries; objects are
restricted to those tokenizing to a single token under the Llama-3.2-3B tokenizer.
Corruptions substitute a different country that predicts a different object
(e.g.\ France/Paris $\to$ Germany/Berlin), preserving sequence length.

We note this dataset is \emph{not} drawn from Meng et al.\ (2022)'s CounterFact set,
as an earlier draft stated. It was authored for this work, and its two relation types
are far narrower than CounterFact's. This limits the cross-task claim in
\S\ref{sec:cross-task} correspondingly: what we test is transfer from IOI to two
geographic relation templates, not to factual recall in general. Factual association
experiments were run on Llama-3.2-3B only. Full dataset:
\texttt{data/factual-assoc/dataset.json}.

% ─────────────────────────────────────────────────────────────────────────────
\section{Results: IOI Circuit Identification}
\label{sec:ioi}

\subsection{Baseline Task Performance}

Both models correctly perform the IOI task across all 100 evaluation examples.
Llama-3.2-3B achieves a mean clean logit difference (LD) of \textbf{5.643}
(SD = 0.773), reflecting a strong preference for the indirect object (IO) token over
the subject (S) token at the final sequence position. Pythia-1.4B achieves a mean
clean LD of \textbf{4.120}, 27\% lower in absolute terms, consistent with its smaller
parameter count and simpler attention architecture. Neither model produces a
negative-LD example: the IO token is ranked above S on every trial in both models,
confirming that the circuit analysis begins from a clean behavioral baseline.

\subsection{IOI Circuit Identification: Llama-3.2-3B}
\label{sec:llama-ioi}

Figure~\ref{fig:heatmap} shows the full activation patching heatmap for Llama-3.2-3B
across all $28 \times 24 = 672$ (layer, head) positions. Two features are immediately
apparent: most heads contribute negligibly to logit-difference recovery under patching,
and a sparse set of heads in the middle-to-late network produces strongly positive
scores.

\begin{figure}[H]
  \centering
  \includegraphics[width=0.92\textwidth]{figures/ioi-patching-heatmap-llama3b}
  \caption{Activation patching heatmap for Llama-3.2-3B ($n = 100$ IOI examples).
    Each cell shows normalized logit-difference recovery when the output of that
    (layer, head) pair is replaced with the corresponding clean-run activation.
    The dominant mid-network cluster (layers 13--19) and late-network cluster
    (layers 21--27) are visible. Data: \texttt{data/ioi/patching-llama3b.json}.}
  \label{fig:heatmap}
\end{figure}

Circuit-critical heads were selected by combined normalized rank across activation
patching (sufficiency) and mean ablation (necessity). Table~\ref{tab:llama-circuit}
lists the top-10 heads by combined rank.

\begin{table}[H]
  \centering
  \caption{Top-10 circuit-critical heads in Llama-3.2-3B by combined normalized rank
    across activation patching (sufficiency) and mean-ablation drop (necessity).
    Relative depth = layer index / $n_\text{layers}$ (zero-indexed).
    Data: \texttt{data/ioi/patching-llama3b.json},
    \texttt{data/ioi/ablation-llama3b.json}.}
  \label{tab:llama-circuit}
  \smallskip
  \begin{tabular}{clccccr}
    \toprule
    Rank & Head    & Patch score & 95\% CI          & Abl.\ drop & \% clean LD & Rel.\ depth \\
    \midrule
    1    & L15H20  & 0.240       & [0.225, 0.255]   & 1.578       & 28.0\%      & 0.536 \\
    2    & L17H17  & 0.098       & [0.089, 0.107]   & 0.929       & 16.5\%      & 0.607 \\
    3    & L13H14  & 0.092       & [0.083, 0.102]   & 0.626       & 11.1\%      & 0.464 \\
    4    & L24H15  & 0.148       & [0.133, 0.164]   & 0.276       &  4.9\%      & 0.857 \\
    5    & L19H1   & 0.099       & [0.094, 0.105]   & 0.443       &  7.8\%      & 0.679 \\
    6    & L21H20  & 0.073       & [0.068, 0.078]   & 0.417       &  7.4\%      & 0.750 \\
    7    & L18H10  & 0.078       & [0.073, 0.084]   & 0.313       &  5.5\%      & 0.643 \\
    8    & L14H0   & 0.111       & [0.099, 0.125]   & 0.017       &  0.3\%      & 0.500 \\
    9    & L27H17  & 0.035       & [0.030, 0.041]   & 0.373       &  6.6\%      & 0.964 \\
    10   & L26H23  & 0.039       & [0.030, 0.047]   & 0.344       &  6.1\%      & 0.929 \\
    \bottomrule
  \end{tabular}
\end{table}

The most striking feature of the Llama-3.2-3B circuit is the dominance of L15H20. Its
ablation drop (1.578) is \textbf{70\% larger} than rank-2 L17H17 (0.929) and
$\mathbf{2.5\times}$ rank-3 L13H14 (0.626). L15H20 also leads the patching ranking at
a score of 0.240, more than 60\% above rank-2 L24H15 (0.148). No other head in the
top-10 approaches this contribution on both metrics simultaneously: L24H15 scores
highly on patching but weakly on ablation (0.276), while L27H17 shows the reverse
pattern. This double-dominance of L15H20---both sufficient and necessary---identifies
it as a hard bottleneck in the IOI circuit at relative depth 0.536.

The circuit spans relative depths 0.46--0.97 with no critical head below 0.46. Three
depth clusters are visible: a mid cluster (0.46--0.54) comprising L13H14, L14H0, and
L15H20; a mid-late cluster (0.60--0.75) comprising L17H17, L18H10, L19H1, and L21H20;
and a late cluster (0.86--0.97) comprising L24H15, L26H23, and L27H17. This
three-cluster structure corresponds qualitatively to the functional zones identified by
Wang et al.\ (2022) in GPT-2 Small, as discussed in Section~\ref{sec:comparison}.

\subsection{IOI Circuit Identification: Pythia-1.4B}
\label{sec:pythia-ioi}

Pythia-1.4B replicates the bottleneck structure of Llama-3.2-3B.
Table~\ref{tab:pythia-circuit} presents the top-10 circuit-critical heads. Note that
the table is ordered by ablation drop rather than combined rank to highlight the
anomalous L1H11 entry discussed below.

\begin{table}[H]
  \centering
  \caption{Top-10 circuit-critical heads in Pythia-1.4B by combined normalized rank
    across activation patching and mean-ablation drop. Table rows are ordered by
    ablation drop (descending) to highlight the L1H11 necessity-without-sufficiency
    anomaly; combined-rank ordering would place L21H3 at position 4 and L12H15 at
    position 5. Data: \texttt{data/ioi/patching-pythia1b.json},
    \texttt{data/ioi/ablation-pythia1b.json},
    \texttt{data/ioi/statistical-validation.json}.}
  \label{tab:pythia-circuit}
  \smallskip
  \begin{tabular}{clccccr}
    \toprule
    Rank & Head   & Patch score & 95\% CI             & Abl.\ drop & \% clean LD & Rel.\ depth \\
    \midrule
    1    & L10H7  & 0.207       & [0.199, 0.216]      & 0.893       & 21.7\%      & 0.417 \\
    2    & L15H15 & 0.155       & [0.134, 0.175]      & 0.551       & 13.4\%      & 0.625 \\
    3    & L22H2  & 0.101       & [0.097, 0.105]      & 0.411       & 10.0\%      & 0.917 \\
    4    & L1H11  & $-$0.003    & [$-$0.011, $+$0.005]& 0.327       &  7.9\%      & 0.042 \\
    5    & L21H3  & 0.056       & [0.052, 0.059]      & 0.237       &  5.8\%      & 0.875 \\
    6    & L17H7  & 0.040       & [0.027, 0.056]      & 0.150       &  3.6\%      & 0.708 \\
    7    & L10H0  & 0.039       & [0.036, 0.041]      & 0.128       &  3.1\%      & 0.417 \\
    8    & L12H15 & 0.051       & [0.048, 0.054]      & 0.124       &  3.0\%      & 0.500 \\
    9    & L16H13 & 0.025       & [0.022, 0.029]      & 0.117       &  2.8\%      & 0.667 \\
    10   & L13H6  & 0.049       & [0.042, 0.056]      & 0.018       &  0.4\%      & 0.542 \\
    \bottomrule
  \end{tabular}
\end{table}

L10H7 leads both rankings: its ablation drop (0.893) is \textbf{62\% larger} than
rank-2 L15H15 (0.551), and its patching score (0.207) is \textbf{34\% larger} than
rank-2 (0.155). The pattern---one head substantially dominant on both sufficiency and
necessity---exactly mirrors the Llama result, though the absolute magnitude of the
ablation drop is smaller (0.893 vs.\ 1.578) consistent with Pythia's lower baseline
LD.

\paragraph{Anomalous early head: L1H11.}
One Pythia-specific finding stands out: L1H11 at relative depth 0.042 (layer 1 of 24)
scores $-0.003$ on activation patching (95\% CI $[-0.011, +0.005]$, straddling zero)
yet produces the fourth-largest ablation drop in the model (0.327, 7.9\% of clean LD).
This dissociation between sufficiency (near zero or slightly negative) and necessity
(substantial) is inconsistent with the role profile of name-mover or S-inhibition
heads, which score positively on both measures. The most parsimonious interpretation
is that L1H11 suppresses interference---for example, dampening an early signal that
would otherwise mislead later heads---rather than directly boosting IO probability.
Ablating it disrupts the clean causal pathway downstream, producing a large LD drop,
but substituting its clean activation into a corrupted context provides no positive
recovery. No Llama-3.2-3B head in the top-10 occupies a comparable depth ($< 0.10$),
suggesting this interference-suppression role may be specific to Pythia's architecture
or training distribution.

\subsection{Comparison with Wang et al.\ (2022): Depth-Zone Conservation}
\label{sec:comparison}

Wang et al.\ (2022) identified three functional head classes in GPT-2 Small at specific
relative depth ranges: duplicate-token/induction heads (depth 0.0--0.42), S-inhibition
heads (0.58--0.67), and name-mover heads (0.75--0.83). Table~\ref{tab:zones} shows
where these zones fall in our two models.

\begin{table}[H]
  \centering
  \caption{Functional depth-zone alignment across GPT-2 Small \citep{wang2022interpretability},
    Llama-3.2-3B, and Pythia-1.4B. Zone boundaries for Llama and Pythia are inferred
    from the depth distribution of critical heads, not directly assigned by functional
    annotation. The name-mover zone boundary is extended from Wang et al.'s original
    0.75--0.83 to 0.75--0.92 to accommodate heads at deeper positions in the modern
    models. Data: \texttt{data/ioi/cross-model-comparison.md}.}
  \label{tab:zones}
  \smallskip
  \small
  \begin{tabular}{p{3.6cm}p{2.4cm}p{3.6cm}p{3.6cm}}
    \toprule
    Zone                    & GPT-2 Small         & Llama-3.2-3B             & Pythia-1.4B \\
    \midrule
    Early (0.00--0.25)      & L0H1, L3H0          & ---                      & L1H11 \\
    Mid-induction (0.40--0.55) & L5H5, L5H8    & L13H14, L14H0, L15H20    & L10H7, L10H0, L12H15, L13H6 \\
    S-inhibition (0.50--0.70) & L7H3--L8H10    & L17H17, L18H10, L19H1    & L15H15, L16H13 \\
    Name-movers (0.75--0.92) & L9H6--L10H0    & L21H20, L24H15, L26H23   & L21H3, L22H2 \\
    \bottomrule
  \end{tabular}
\end{table}

\begin{figure}[H]
  \centering
  \includegraphics[width=0.85\textwidth]{figures/ioi-circuit-diagram}
  \caption{Schematic circuit diagram for the IOI circuit across GPT-2 Small,
    Llama-3.2-3B, and Pythia-1.4B, annotated by functional zone.
    Data: \texttt{data/ioi/cross-model-comparison.md}.}
  \label{fig:circuit-diagram}
\end{figure}

All three functional zones from Wang et al.\ are present in both modern models at
compatible relative depths. The mid-induction zone (0.40--0.55) is the densest cluster
in all three models; the name-mover zone (0.75--0.92) contains the heads with highest
late-network patching scores. One notable divergence: the very-early zone (0.00--0.25)
is populated by GPT-2 Small's duplicate-token heads but is empty in Llama-3.2-3B's
top-10 and represented only by the anomalous L1H11 in Pythia.

A second divergence is the \textbf{bottleneck vs.\ distributed} structure. Wang et
al.\ found 26 circuit components with moderate individual contributions; in both Llama
and Pythia, one head accounts for 22--28\% of clean LD on its own, substantially above
any single head's contribution in GPT-2 Small. This could reflect scale effects (larger
models can concentrate circuit function in individual heads), differences in our
selection methodology, or task-distribution differences in our IOI dataset. We return
to this in Section~\ref{sec:discussion}.

\subsection{Statistical Significance and Patching Verification}

Bootstrap confidence intervals (1{,}000 resamples, 95\% level, seed 42) were computed
for all heads with mean patching score $\geq 0.030$. For the top-ranked heads in each
model, the effect sizes are large relative to uncertainty:

\begin{itemize}
  \item \textbf{L15H20 (Llama):} mean 0.240, 95\% CI [0.225, 0.255], CI width 0.030.
    The lower bound (0.225) exceeds the rank-3 head's mean (0.092) by $2.4\times$.
  \item \textbf{L10H7 (Pythia):} mean 0.207, 95\% CI [0.199, 0.216], CI width 0.017.
    The lower CI bound (0.199) exceeds rank-2 (0.155) by 28\%.
\end{itemize}

The margin varies substantially across the top-10, and we state it precisely rather
than as a blanket claim. Ranking heads by the ratio of CI lower bound to CI width:
seven of ten Llama heads and six of ten Pythia heads exceed $3\times$. The remainder
do not---L27H17 ($2.7\times$) and L26H23 ($1.8\times$) in Llama, and L17H7
($0.9\times$) in Pythia, whose interval $[0.027, 0.056]$ is wide relative to its mean.
These four heads (including L1H11, below) are the weakest members of their circuits and
should be read as tentative. An earlier draft asserted the $3\times$ margin for all
top-10 heads in both models; that was wrong, and contradicted by the tables in this
same section.

\paragraph{No correction for multiple comparisons.} The bootstrap intervals above are
computed on the same 100 examples used to select the heads, after screening 672
(layer, head) pairs in Llama and 384 in Pythia. They are therefore conditional on
selection and are optimistic: the winner's-curse bias inflates the top-ranked means,
and no family-wise or FDR correction is applied. The three or four dominant heads in
each model have effects far too large for this to matter---L15H20's lower bound
exceeds the rank-3 mean by $2.4\times$---but for heads in the 0.02--0.05 score range
the intervals should not be read as calibrated. A clean fix is available and cheap:
select heads on one half of the example set and estimate effects on the other.

The Pythia L1H11 anomaly is confirmed: its patching
CI $[-0.011, +0.005]$ straddles zero (the only such case in either model's top-10),
whereas its ablation drop (0.327) is large and unambiguous. This dissociation is the
empirical signature of necessity without sufficiency, and constitutes a primary
verification in our results: by simultaneously measuring what happens when a head's
activation is \emph{replaced} (patching; sufficiency) versus \emph{removed} (ablation;
necessity), we detect heads whose role is suppressive rather than generative---a
distinction that patching-only or ablation-only protocols would miss.

% ─────────────────────────────────────────────────────────────────────────────
\section{Results: Cross-Architecture and Cross-Task Generalization}
\label{sec:generalization}

\subsection{Cross-Architecture Generalization}

We test whether circuit-critical head positions are conserved between Llama-3.2-3B and
Pythia-1.4B by matching heads greedily by minimum relative-depth distance, with a
$\pm$0.075 tolerance. Table~\ref{tab:cross-arch} reports all matched pairs.

\begin{figure}[H]
  \centering
  \includegraphics[width=0.80\textwidth]{figures/ioi-cross-model-comparison}
  \caption{Cross-model comparison of circuit-critical head positions by relative depth
    for Llama-3.2-3B and Pythia-1.4B. Matched pairs are connected by dashed lines;
    unmatched (model-specific) heads are shown without connections.
    Data: \texttt{data/ioi/cross-model-comparison.md}.}
  \label{fig:cross-model}
\end{figure}

\begin{table}[H]
  \centering
  \caption{Cross-architecture matched head positions (Llama-3.2-3B vs.\ Pythia-1.4B,
    tolerance $\pm$0.075). The $\sim$0.54 pairing (L15H20 $\leftrightarrow$ L13H6)
    reflects depth-position conservation but not functional equivalence: L15H20 is
    Llama's dominant bottleneck head (ablation drop 1.578, 28.0\% of clean LD) while
    L13H6 is Pythia's rank-10 entry (ablation drop 0.018, 0.4\% of clean LD).
    Data: \texttt{data/ioi/cross-model-comparison.md}.}
  \label{tab:cross-arch}
  \smallskip
  \begin{tabular}{cccccc}
    \toprule
    Zone   & Llama head & Llama depth & Pythia head & Pythia depth & $|\Delta|$ \\
    \midrule
    $\sim$0.50 & L14H0  & 0.500 & L12H15 & 0.500 & 0.000 \\
    $\sim$0.54 & L15H20 & 0.536 & L13H6  & 0.542 & 0.006 \\
    $\sim$0.68 & L19H1  & 0.679 & L16H13 & 0.667 & 0.012 \\
    $\sim$0.92 & L26H23 & 0.929 & L22H2  & 0.917 & 0.012 \\
    $\sim$0.62 & L17H17 & 0.607 & L15H15 & 0.625 & 0.018 \\
    $\sim$0.86 & L24H15 & 0.857 & L21H3  & 0.875 & 0.018 \\
    $\sim$0.73 & L21H20 & 0.750 & L17H7  & 0.708 & 0.042 \\
    $\sim$0.46 & L13H14 & 0.464 & L10H7  & 0.417 & 0.048 \\
    \bottomrule
  \end{tabular}
\end{table}

\textbf{8 of 10 circuit-critical head positions match within $\pm$0.075 relative
depth.} The two model-specific heads are: Llama's L18H10 (depth 0.643, no Pythia
match) and L27H17 (depth 0.964, near-final layer, consistent with Llama's greater
depth permitting an extra output-adjustment head); and Pythia's L1H11 (depth 0.042,
discussed above) and L10H0 (depth 0.417, co-located with L10H7 in the same layer,
forming a two-head cluster not observed in Llama).

\subsubsection{The 80\% match rate is not distinguishable from chance}
\label{sec:depth-null}

Taken alone, ``8 of 10'' invites the reading that circuit positions are conserved
across architectures. That reading does not survive a null test, and we report the
null here rather than in an appendix because the claim is the one most likely to be
carried away from this paper.

We ask how often the same greedy matcher, with the same $\pm$0.075 tolerance, would
pair 8 or more of 10 head positions if the heads were placed at random. Two nulls are
reported: a \emph{uniform} null drawing 10 layers uniformly from each model's full
stack, and an \emph{informed} null drawing only from the layer band where
circuit-critical heads actually sit ($d \geq 0.40$ in both models). The informed null
is the fairer comparison, because the claim under test is that specific depths
correspond---not that circuits live in the second half of the network, which both
models' results establish independently.

\begin{table}[H]
  \centering
  \caption{Permutation null for cross-architecture depth matching (20{,}000 draws each,
    seed 42). The observed 8/10 is within the range routinely produced by random head
    placement. Data: \texttt{data/sae-analysis/corrections.json}.}
  \label{tab:depth-null}
  \smallskip
  \begin{tabular}{lccc}
    \toprule
    & Mean matches & Match rate & $\Pr(\geq 8/10)$ \\
    \midrule
    Observed                       & 8.00 & 80.0\% & --- \\
    Uniform null (all layers)      & 5.88 & 58.8\% & \textbf{0.098} \\
    Informed null ($d \geq 0.40$)  & 6.87 & 68.7\% & \textbf{0.32} \\
    \bottomrule
  \end{tabular}
\end{table}

Neither null rejects at any conventional level. Under the informed null---random heads
drawn from the band where real circuit heads live---random sets match at 69\% on
average, and reach 80\% or better about a third of the time.

The reason is a mismatch of scale between the tolerance and the data. A $\pm$0.075
window on relative depth spans roughly two layers in Llama (of 28) and two in Pythia
(of 24). Ten heads per model, drawn from a band of about 17 usable layers, cannot
avoid landing within two layers of one another; the greedy matcher then takes the
closest available pairing. The tolerance was chosen to be narrower than the spacing
between Wang et al.'s functional zones, which is a reasonable criterion for
\emph{zone assignment} but is far too coarse to establish \emph{position
correspondence} between individual heads.

\paragraph{The matches are also functionally uninformative.} Table~\ref{tab:cross-arch}
pairs Llama's L15H20---the dominant bottleneck, ablation drop 1.578---with Pythia's
L13H6, whose ablation drop is 0.018, two orders of magnitude smaller. Meanwhile
Pythia's own dominant head, L10H7, is matched to Llama's rank-3 L13H14. The procedure
never had access to head importance, so nothing constrains it to pair heads playing
similar roles, and it does not. Any conserved structure that exists here is not
captured by depth proximity alone.

\paragraph{What would establish depth conservation.} A valid test needs at least three
things this analysis lacks: functional role assignment for each head (via path
patching and attention-pattern analysis) so that matching is role-aware rather than
purely positional; a tolerance derived from the empirical precision of depth
estimates rather than from zone spacing; and more than two models, so that the
question becomes whether a depth ordering replicates across a population rather than
whether two specific stacks happen to align. We regard the three-cluster structure
visible in both models (\S\ref{sec:llama-ioi}, \S\ref{sec:pythia-ioi}) as suggestive
and the quantitative conservation claim as unsupported.

The matched pairs exhibit a three-cluster depth structure in both models:

\begin{itemize}
  \item \textbf{Mid cluster (0.40--0.55):} Three Llama heads, four Pythia heads.
    Pythia is denser here, consistent with its smaller total layer count producing more
    critical heads in the early-to-mid range.
  \item \textbf{Mid-late cluster (0.60--0.75):} Four Llama heads, three Pythia heads.
  \item \textbf{Late cluster (0.85--0.97):} Three Llama heads, two Pythia heads. Llama
    is denser here, reflecting its two additional layers (28 vs.\ 24) creating room
    for extra late-network mechanisms.
\end{itemize}

The depth-difference distribution across matched pairs is tightly concentrated: six of
eight pairs differ by $\leq 0.02$ in relative depth, and the maximum difference is
0.048. Considering that the tolerance was set at 0.075, the actual conservation is
substantially tighter than the threshold requires.

\subsection{Path Patching: Edge Structure Around the Bottleneck}
\label{sec:path-results}

Table~\ref{tab:path} reports all ten ordered edges among Llama's five highest-scoring
heads. The picture is consistent with L15H20 acting as an early bottleneck that is fed
from below and that feeds the late-network heads only weakly.

\begin{table}[H]
  \centering
  \caption{Edge-level path patching in Llama-3.2-3B ($n = 100$). Score is normalized
    logit-difference recovery attributable to the direct sender$\to$receiver edge.
    Positive means the edge carries information that supports correct IO prediction;
    negative means patching the edge \emph{increases} logit difference above the clean
    baseline. Data: \texttt{data/ioi/path-patching-llama3b-real.json}.}
  \label{tab:path}
  \smallskip
  \begin{tabular}{llrrl}
    \toprule
    Sender & Receiver & Score & $\Delta$ LD & Interpretation \\
    \midrule
    L14H0  & L15H20 & $+0.165$ & $-1.879$ & Dominant edge \\
    L15H20 & L19H1  & $+0.038$ & $-0.431$ & Weak forward path \\
    L15H20 & L24H15 & $+0.037$ & $-0.420$ & Weak forward path \\
    L17H17 & L19H1  & $+0.021$ & $-0.241$ & \\
    L14H0  & L24H15 & $+0.017$ & $-0.195$ & \\
    L19H1  & L24H15 & $+0.016$ & $-0.178$ & \\
    L17H17 & L24H15 & $+0.015$ & $-0.169$ & \\
    L14H0  & L19H1  & $+0.014$ & $-0.154$ & \\
    L14H0  & L17H17 & $-0.056$ & $+0.627$ & Suppressive \\
    L15H20 & L17H17 & $-0.113$ & $+1.273$ & Suppressive \\
    \bottomrule
  \end{tabular}
\end{table}

Three observations follow.

\paragraph{L14H0 $\to$ L15H20 is the dominant edge.} At $0.165$ it is $4.4\times$ the
next-largest positive edge. This resolves the anomaly noted in
\S\ref{sec:llama-ioi}: L14H0 has a strong patching score ($0.111$, rank 3) but a
negligible ablation drop ($0.017$, $0.3\%$ of clean LD). A head can matter greatly as a
\emph{source} for a downstream computation while contributing almost nothing when
removed in isolation, if the downstream head can compensate from other inputs. L14H0
looks like a feeder into the bottleneck rather than an independent contributor.

\paragraph{L15H20's forward edges are weak.} Its contributions to L19H1 ($0.038$) and
L24H15 ($0.037$) are an order of magnitude below the edge feeding it. Combined with its
large ablation drop ($1.578$), this suggests L15H20 writes substantially to the final
position directly, rather than acting through the later heads we measured. We did not
patch the L15H20~$\to$~logit edge itself, so this is an inference from the absence of
strong mediated paths and not a direct measurement.

\paragraph{Two edges are suppressive.} Patching L15H20~$\to$~L17H17 \emph{raises} logit
difference by $1.273$ above the clean baseline (score $-0.113$), and
L14H0~$\to$~L17H17 does likewise. The natural reading is that these edges carry a
signal that normally damps the IO logit, and that replacing it with the corrupted-run
value removes the damping. Since L17H17 is rank 2 by ablation ($0.929$), this is not a
peripheral component: the circuit appears to contain both excitatory and inhibitory
paths into the same head. Wang et al.\ report an analogous negative-name-mover
phenomenon in GPT-2 Small, and this is the closest correspondence to their functional
taxonomy that our data directly supports.

\paragraph{Scope.} Ten edges among five heads is a small fraction of the circuit, and
these scores carry no confidence intervals---the path-patching run recorded per-edge
means only. Pythia was not measured at all. We report the structure as a directional
finding, not a validated circuit diagram.

\subsection{Cross-Task Transfer: Factual Association}
\label{sec:cross-task}

To test whether the circuit structure identified for IOI transfers to a related but
distinct task, we applied the same activation patching protocol to factual association
prompts ($n = 50$, Llama-3.2-3B only). Prompts had the form ``The capital of [country]
is \_\_\_'' with the country name corrupted by substitution to a different country; the
patching score measures each head's causal contribution to correct capital prediction.
Table~\ref{tab:factual-assoc} reports the top-10 heads.

\begin{table}[H]
  \centering
  \caption{Top-10 heads by activation patching score for factual association
    (Llama-3.2-3B, $n = 50$). The ``Matching IOI layer?'' column indicates whether
    the head's layer appears among the IOI top-9 heads.
    Data: \texttt{data/factual-assoc/patching-llama3b.json}.}
  \label{tab:factual-assoc}
  \smallskip
  \begin{tabular}{clccl}
    \toprule
    Rank & Head   & Patch score & Rel.\ depth & Matching IOI layer? \\
    \midrule
    1  & L15H17 & 0.420 & 0.536 & \textbf{Yes} (IOI rank 1: L15H20) \\
    2  & L21H2  & 0.418 & 0.750 & \textbf{Yes} (IOI rank 6: L21H20) \\
    3  & L27H5  & 0.105 & 0.964 & \textbf{Yes} (IOI rank 9: L27H17) \\
    4  & L17H18 & 0.088 & 0.607 & \textbf{Yes} (IOI rank 2: L17H17) \\
    5  & L13H18 & 0.052 & 0.464 & \textbf{Yes} (IOI rank 3: L13H14) \\
    6  & L25H8  & 0.045 & 0.893 & No \\
    7  & L26H20 & 0.035 & 0.929 & No \\
    8  & L13H19 & 0.026 & 0.464 & Yes (layer 13) \\
    9  & L21H0  & 0.025 & 0.750 & Yes (layer 21) \\
    10 & L26H19 & 0.019 & 0.929 & No \\
    \bottomrule
  \end{tabular}
\end{table}

The top-5 factual association heads occupy the same five layers as five of the top-9
IOI heads (layers 13, 15, 17, 21, and 27), but activate \emph{different heads within
those layers}: IOI uses L15H20 while FA uses L15H17; IOI uses L21H20 while FA uses
L21H2; IOI uses L13H14 while FA uses L13H18. This pattern is consistent with
\textbf{layer-level circuit topology conservation across tasks} with \textbf{head-level
specialization by task}---the same network locations implement different input-output
functions depending on which specific circuit the task activates.

Two additional features of the factual association results are noteworthy. First, the
top-2 FA heads (0.420, 0.418) have substantially higher patching scores than the IOI
rank-1 head (0.240), and the two scores are nearly equal to each other, suggesting FA
uses a more concentrated two-head bottleneck while IOI distributes attribution more
broadly across the mid-to-late network. Second, layer 15 is the single most critical
layer for both tasks---the dominant head in each circuit (L15H20 for IOI, L15H17 for
FA) resides at the same relative depth (0.536), which may reflect layer 15's position
at the transition between Llama's induction-like and output-writing computational
stages.

% ─────────────────────────────────────────────────────────────────────────────
\section{Efficiency Analysis}
\label{sec:efficiency}

The activation-patching and mean-ablation protocols in this paper are computationally
intensive by design: a complete IOI sweep requires 674 forward passes for Llama-3.2-3B
and 386 passes for Pythia-1.4B, all batched over 100 examples. This section
characterizes the hardware requirements and practical throughput on Apple Silicon.

\paragraph{Memory footprint.}
Llama-3.2-3B in bfloat16 requires approximately 6.4\,GB for model weights alone. The
activation cache adds $n_\text{examples} \times n_\text{layers} \times n_\text{heads}
\times L_\text{seq} \times d_\text{head} \times 2$ bytes per forward-pass
cache---roughly 0.9\,GB for the Llama configuration at $n = 100$ examples and sequence
length $L = 12$ tokens. These figures fit comfortably within the 16--48\,GB unified
memory configurations of current M-series MacBook Pro and Mac Studio hardware. The
model weights and activation cache share the same physical memory as the CPU address
space, eliminating any PCIe transfer bottleneck that would be present in a separate
discrete GPU configuration.

\paragraph{Throughput.}
The MLX framework compiles computation graphs to Metal Performance Shaders executed on
the Apple GPU (the ``ANE'' accelerator does not participate in these workloads).
Preliminary timing measurements indicate that the full Llama-3.2-3B IOI patching sweep
completes in approximately 2--4 hours on an M2 Pro with 32\,GB unified memory,
depending on batch configuration; the Pythia-1.4B sweep completes in under 2 hours.
Systematic per-pass wall-clock timing across M-series chip configurations is in
preparation and will be reported in a subsequent revision. The ablation sweep is
comparable in cost to the patching sweep ($1 + n_L \times n_H$ vs.\
$2 + n_L \times n_H$ passes).

\paragraph{Reproducibility.}
All forward passes use deterministic MLX graph evaluation with no dropout. Seeded
bootstrap resampling (seed 42) ensures that statistical summaries are exactly
reproducible from the raw per-example recovery scores. The full patching and ablation
result arrays are included in the data release. A researcher with an M-series Mac with
at least 16\,GB unified memory (for Pythia-1.4B) or 24\,GB (for Llama-3.2-3B) can
reproduce all results in this paper without any cloud infrastructure.

% ─────────────────────────────────────────────────────────────────────────────
\section{Discussion and Conclusion}
\label{sec:discussion}

\paragraph{Three-zone structure without depth conservation.}
Both Llama-3.2-3B and Pythia-1.4B exhibit a three-zone qualitative structure---early
token-marking heads, mid-network suppression, and late name-mover heads---broadly
compatible with the zones Wang et al.\ (2022) identify in GPT-2 Small. This
structural parallel holds despite a 3--7$\times$ difference in parameter count, the
introduction of grouped-query attention and rotary position embeddings in Llama, and
substantially different pretraining distributions. In that descriptive sense, both
models appear to solve IOI through a qualitatively similar functional decomposition.

The quantitative form of this claim---that head positions are metrically conserved by
relative depth---does not hold. Matching circuit-critical heads between the two models
at $\pm$0.075 relative depth yields 8 of 10 pairs, but randomly sampled head sets
achieve 59--69\% under the same procedure ($p = 0.098$ and $p = 0.32$ under uniform
and band-restricted nulls; \S\ref{sec:depth-null}). The 80\% figure is not
distinguishable from chance, and the best-matched pair links Llama's dominant head
(L15H20, ablation drop 1.578) to Pythia's weakest circuit member (L13H6, ablation
drop 0.018)---two orders of magnitude apart. We retain the three-cluster observation
as descriptive and withdraw the depth-conservation claim, consistent with the abstract.

\paragraph{Bottleneck structure at scale.}
A striking departure from Wang et al.\ is the degree of functional concentration in
our models: L15H20 in Llama accounts for 28.0\% of clean logit difference on its own,
and L10H7 in Pythia accounts for 21.7\%. GPT-2 Small's 26-component circuit
distributes attribution more evenly, with no single head dominating to this degree.
Several hypotheses can explain this. Scale may allow models to assign more of a
computation to a single highly specialized head, since a larger total head count
provides enough redundancy elsewhere in the network. Alternatively, our selection
criteria may be stricter than Wang et al.'s, causing us to exclude peripheral
contributors that they retained. A third possibility we cannot exclude is that the
comparison is not like-for-like: Wang et al.\ measure a circuit assembled and
validated as a set, whereas our top-10 lists are ranked marginal effects, and the two
need not decompose the same way.

The path-patching results in \S\ref{sec:path-results} bear on the first question for
Llama. L15H20's forward edges into the later heads we measured (0.038 and 0.037) are
an order of magnitude weaker than the L14H0~$\to$~L15H20 edge feeding it (0.165),
which is consistent with L15H20 writing to the output directly rather than through
downstream mediation. We did not patch the head-to-logit edge itself, and Pythia was
not path-patched at all, so this is suggestive rather than settled.

\paragraph{Depth-position conservation does not imply functional equivalence.}
The cross-model comparison pairs L15H20 (Llama) with L13H6 (Pythia) at relative depth
$\sim$0.54. These heads share a depth zone but are not functionally equivalent:
L15H20's ablation drop (1.578) exceeds L13H6's (0.018) by two orders of magnitude.
The depth-position match reflects that \emph{something important} happens in both
models near relative depth 0.54---it does not imply that the two heads play the same
role. L15H20 is Llama's dominant bottleneck; L13H6 is on the periphery of Pythia's
circuit. Future work assigning fine-grained functional roles (via path patching and
attention-pattern analysis) will clarify whether this zone hosts the S-inhibition or
induction-copying function in each model.

\paragraph{Layer-level topology and task transfer.}
The factual association results indicate that Llama-3.2-3B has a stable set of
\emph{critical layers}---particularly layers 13, 15, 17, 21, and 27---that appear
across tasks with different head assignments within each layer. This is consistent with
the hypothesis that certain layers develop as ``computation hubs'' during pretraining,
acquiring dense connectivity patterns that many downstream circuits route through,
while head-level specialization allows different circuits to use different heads within
the same hub. Layer 15 is particularly notable: the rank-1 head for both IOI
(L15H20, patch score 0.240) and factual association (L15H17, patch score 0.420) resides
here, suggesting it occupies a structural bottleneck position in Llama's residual
stream.

\paragraph{SAE-based feature attribution as a next step.}
The circuit-level results reported here identify \emph{which} heads matter; they do not
explain \emph{why} those heads matter. A natural next step is to apply sparse
autoencoder (SAE) decomposition to the residual stream at the positions that L15H20
and L10H7 read from, and to the outputs they write back. This would reveal which
monosemantic features drive the dominant heads' attention patterns and output
projections. If L15H20 reads a ``name at position $k$'' feature and writes a
``preferred next token'' feature, the interpretive picture is complete; if the readout
is more distributed across multiple SAE features, the circuit analysis is incomplete
without the feature-level complement. We plan to pursue this in a companion paper
(cf.\ \citealt{elhage2022toy}).

\paragraph{Implications for consumer-hardware mechanistic interpretability.}
All experiments in this paper were conducted on Apple Silicon without access to NVIDIA
GPU clusters. The results demonstrate that full activation-patching circuit analysis of
1--3B parameter models is computationally tractable on modern M-series hardware,
requiring approximately 2--4 hours of wall time per model for the full IOI sweep. The
key enabling factor is Apple's unified memory architecture: a 24--48\,GB unified
memory configuration can hold both model weights and activation caches without
offloading, making the per-pass latency competitive with small GPU cluster runs. For
models beyond 7B parameters, the memory constraint becomes binding on current
hardware---a limitation that may be partially addressed by mixed-precision loading or
activation recomputation at the cost of higher pass count. We hope that releasing the
SwiftSci Interp tooling lowers the barrier for researchers to replicate and extend
circuit analyses without requiring institutional compute resources.

\paragraph{Limitations.}

\textbf{Circuit faithfulness collapses on held-out data.} Wang et al.\ validate their
circuit as a \emph{set}, reporting that it recovers 87\% of clean performance when
everything outside it is ablated. We ran the same evaluation. On the original
100-example single-template dataset (v1), Llama's circuit-only logit difference
recovers $F_\text{faithfulness} = 0.760$ and Pythia's recovers $0.854$---both
competitive with Wang et al.\ On the held-out 200-example dataset (v2, expanded name
set), faithfulness collapses: Llama drops to $F_\text{faithfulness} = 0.228$ and
Pythia to $-0.028$. Pythia's circuit-only model scores \emph{below} the
fully-ablated baseline on v2, meaning the ten heads selected on v1 are collectively
anti-predictive on out-of-distribution names. Joint ablation of all ten Llama heads on
v2 finds $F_\text{complete} = 0.213$, confirming those heads are necessary (removing
them leaves the model at 21\% of clean performance) even when the circuit alone is
insufficient; the leave-one-in analysis identifies L17H17 as the largest single
contributor ($\Delta F = +0.173$). Per-head marginal patching scores are still not
additive (our Llama top-10 sum to 1.013), and faithfulness of individual heads does
not imply faithfulness of the set on a different distribution. The v1$\to$v2 collapse
is a direct consequence of the single-template design flagged above.

\textbf{Single template, single name ordering.} All 100 IOI examples instantiate one
sentence frame, ``When \{S\} and \{IO\} went to the store, \{S\} gave a bottle to'',
over 20 names. Wang et al.\ use 15 templates and balance both name orderings (ABBA and
BABA) precisely because a single frame confounds circuit structure with positional
structure. Heads we identify as circuit-critical may be responding to a fixed token
position rather than to the IOI computation. This limitation applies to every result
in the paper and is the first thing a replication should change.

\textbf{Selection and inference on the same data.} Heads were chosen by ranking 672
(Llama) and 384 (Pythia) candidates, and their confidence intervals were then
bootstrapped on the same 100 examples, with no multiple-comparison correction
(\S\ref{sec:stats}). The dominant heads survive any reasonable correction; the
0.02--0.05 band does not obviously.

\textbf{Depth conservation is unsupported.} The 80\% cross-architecture match rate
does not exceed a permutation null (\S\ref{sec:depth-null}), and the matching pairs
functionally dissimilar heads. We retain the three-cluster observation as descriptive
and withdraw the quantitative claim.

\textbf{Narrow factual-association probe.} The cross-task experiment uses 50
hand-built examples across two geographic relation templates on one model---not
CounterFact, as an earlier draft stated. The layer-overlap statistic also lacks a
null: both tasks' heads concentrate in layers 13--27, so some overlap is expected by
construction. Replication on Pythia and on genuinely diverse relations is needed.

\textbf{Partial path patching.} Ten edges among five Llama heads, no confidence
intervals, and no Pythia measurement (\S\ref{sec:path-results}). The
``hard bottleneck'' characterization of L15H20 remains provisional.

\textbf{Unverified functional roles.} The interpretation of L1H11 (Pythia) as an
interference-suppression head is consistent with its patching/ablation dissociation
but has not been checked against attention patterns. More broadly, we assign no head
to Wang et al.'s functional classes on direct evidence; the zone assignments in
Table~\ref{tab:zones} are inferred from depth alone.

\textbf{Efficiency figures are estimates.} The wall-clock numbers in
\S\ref{sec:efficiency} are preliminary readings from a single machine, not a
benchmark. Systematic per-pass timing across M-series configurations is forthcoming.
This is a real gap in a paper whose stated contribution includes tractability on
Apple Silicon.

\paragraph{Conclusion.}
We set out to test whether the IOI circuit Wang et al.\ (2022) found in GPT-2 Small
generalizes to modern architectures. The answer we can support is narrower than the
one we initially drew.

What holds: both Llama-3.2-3B and Pythia-1.4B solve IOI with a sparse, strongly
bottlenecked circuit. A single head in each (L15H20, L10H7) accounts for 22--28\% of
clean logit difference on its own, well above any single head's share of the
distributed 26-component GPT-2 Small circuit. Critical heads in both models fall into
three depth clusters broadly compatible with Wang et al.'s functional zones. Path
patching in Llama shows L15H20 fed dominantly by L14H0, weakly connected forward, and
receiving suppressive input paths---a shape consistent with direct output writing.

What does not hold: the quantitative depth-conservation claim. Matching head positions
across the two models at $\pm$0.075 relative depth succeeds for 8 of 10 heads, but
random head sets achieve 59--69\% under the same procedure, and the observed rate is
not significant under either null we tested ($p = 0.098$, $p = 0.32$). The matching
also pairs Llama's dominant head with Pythia's weakest, so it carries no functional
information. We report this because ``80\% of positions conserved'' is a memorable
number, and it should not be carried forward.

The largest gap revealed by the circuit-level evaluation is generalization: circuit
faithfulness drops from 0.760 (Llama) and 0.854 (Pythia) on the original
100-example dataset to 0.228 and $-0.028$ respectively on a held-out 200-example
set. The circuits are real on the training distribution---joint ablation confirms
they account for the bulk of IOI performance---but they do not describe the models'
general IOI behaviour. Combined with a single-template dataset, that leaves the
central question---whether these are the same circuit as Wang et al.'s or merely
circuits in similar places---open.

All experiments ran on Apple Silicon using MLX, and the SwiftSci Interp harness is
released alongside this paper to make circuit analysis accessible without GPU cluster
infrastructure. That part of the contribution we regard as demonstrated.

% ─────────────────────────────────────────────────────────────────────────────
\section*{Acknowledgments}

The author thanks the SwiftSci community for support with the MLX activation-patching
infrastructure, and the mechanistic interpretability research community for making
foundational datasets, models, and tools publicly available. The IOI dataset
construction followed the protocol of Wang et al.\ (2022); the CounterFact evaluation
set used in the factual association experiments was constructed by Meng et al.\ (2022).
Pythia models are provided by EleutherAI under the Apache 2.0 license; Llama-3.2-3B
is provided by Meta under the Llama 3.2 Community License Agreement.

% ─────────────────────────────────────────────────────────────────────────────
\section*{Data Availability}

All experimental datasets, intermediate patching and ablation result arrays, and
bootstrap confidence interval files referenced in this paper are available in the
project data repository. The SwiftSci Interp activation-patching library, including
the \texttt{PatchableAttention} module for Llama and Pythia and the MLX-compatible
patching sweep harness, is available at:
\begin{center}
  \url{https://github.com/jmelton-research/swiftsci-interp}
\end{center}
Figure source files are in \texttt{figures/circuit-tracing/}. Raw data files are in
\texttt{data/ioi/} and \texttt{data/factual-assoc/}. All data was generated with seed
42 and is exactly reproducible using the released code.

% ─────────────────────────────────────────────────────────────────────────────
\bibliographystyle{plainnat}
\begin{thebibliography}{99}

\bibitem[Ainslie et al.(2023)]{ainslie2023gqa}
Ainslie, J., Lee-Thorp, J., de Jong, M., Zemlyanskiy, Y., Lebrón, F., \& Sanghai, S.
(2023).
GQA: Training Generalized Multi-Query Transformer Models from Multi-Head Checkpoints.
In \emph{Proceedings of the 2023 Conference on Empirical Methods in Natural Language
Processing (EMNLP)}.

\bibitem[Biderman et al.(2023)]{biderman2023pythia}
Biderman, S., Schoelkopf, H., Anthony, Q., Bradley, H., O'Brien, K., Hallahan, E.,
Khan, M.~A., Goh, S., Galarraga, H.~W., Phang, J., Presser, C., Purohit, H.,
Reynolds, L., Tow, J., Wang, B., \& Weinbach, S.
(2023).
Pythia: A Suite for Analyzing Large Language Models Across Training and Scaling.
In \emph{Proceedings of the 40th International Conference on Machine Learning (ICML)}.

\bibitem[Elhage et al.(2021)]{elhage2021mathematical}
Elhage, N., Nanda, N., Olsson, C., Henighan, T., Joseph, N., Mann, B., Askell, A.,
Bai, Y., Chen, A., Conerly, T., DasSarma, N., Drain, D., Ganguli, D., Hatfield-Dodds,
Z., Hernandez, D., Jones, A., Kernion, J., Lovitt, L., Ndousse, K., Amodei, D.,
Brown, T., Clark, J., Kaplan, J., McCandlish, S., \& Olah, C.
(2021).
A Mathematical Framework for Transformer Circuits.
\emph{Transformer Circuits Thread}.
\url{https://transformer-circuits.pub/2021/framework/index.html}

\bibitem[Elhage et al.(2022)]{elhage2022toy}
Elhage, N., Hume, T., Olsson, C., Schiefer, N., Henighan, T., Kravec, S., Hatfield-Dodds,
Z., Lasenby, R., Drain, D., Chen, C., Grosse, R., McCandlish, S., Kaplan, J.,
Amodei, D., Wattenberg, M., \& Olah, C.
(2022).
Toy Models of Superposition.
\emph{Transformer Circuits Thread}.
\url{https://transformer-circuits.pub/2022/toy_model/index.html}

\bibitem[Goldowsky-Dill et al.(2023)]{goldowsky2023localizing}
Goldowsky-Dill, N., MacLeod, C., Shlegeris, L., \& Hatte, N.
(2023).
Localizing Model Behavior with Path Patching.
\emph{arXiv preprint arXiv:2304.05969}.

\bibitem[Meng et al.(2022)]{meng2022locating}
Meng, K., Bau, D., Andonian, A., \& Belinkov, Y.
(2022).
Locating and Editing Factual Associations in GPT.
In \emph{Advances in Neural Information Processing Systems (NeurIPS)}, Vol.~35.

\bibitem[Meta(2024)]{meta2024llama3}
Meta.
(2024).
The Llama 3 Herd of Models.
\emph{arXiv preprint arXiv:2407.21783}.

\bibitem[Nanda \& Lawrence(2022)]{nanda2022transformerlens}
Nanda, N., \& Lawrence, J.
(2022).
TransformerLens.
\url{https://github.com/neelnanda-io/TransformerLens}

\bibitem[Su et al.(2021)]{su2021roformer}
Su, J., Lu, Y., Pan, S., Murtadha, A., Wen, B., \& Liu, Y.
(2021).
RoFormer: Enhanced Transformer with Rotary Position Embedding.
\emph{arXiv preprint arXiv:2104.09864}.

\bibitem[Conmy et al.(2023)]{conmy2023acdc}
Conmy, A., Mavor-Parker, A.~N., Lynch, A., Heimersheim, S., \& Garriga-Alonso, A.
(2023). Towards Automated Circuit Discovery for Mechanistic Interpretability.
In \emph{Advances in Neural Information Processing Systems (NeurIPS)}, Vol.~36.

\bibitem[Hanna et al.(2023)]{hanna2023greaterthan}
Hanna, M., Liu, O., \& Variengien, A. (2023). How does GPT-2 compute greater-than?
Interpreting mathematical abilities in a pre-trained language model.
In \emph{Advances in Neural Information Processing Systems (NeurIPS)}, Vol.~36.

\bibitem[Heimersheim \& Nanda(2024)]{heimersheim2024patching}
Heimersheim, S., \& Nanda, N. (2024). How to use and interpret activation patching.
arXiv preprint arXiv:2404.15255.

\bibitem[Lieberum et al.(2023)]{lieberum2023scale}
Lieberum, T., Rahtz, M., Kram\'ar, J., Nanda, N., Irving, G., Shah, R., \&
Mikulik, V. (2023). Does Circuit Analysis Interpretability Scale? Evidence from
Multiple Choice Capabilities in Chinchilla. arXiv preprint arXiv:2307.09458.

\bibitem[Merullo et al.(2024)]{merullo2024reuse}
Merullo, J., Eickhoff, C., \& Pavlick, E. (2024). Circuit Component Reuse Across
Tasks in Transformer Language Models. In \emph{International Conference on Learning
Representations (ICLR)}.

\bibitem[Zhang \& Nanda(2024)]{zhang2024bestpractices}
Zhang, F., \& Nanda, N. (2024). Towards Best Practices of Activation Patching in
Language Models: Metrics and Methods. In \emph{International Conference on Learning
Representations (ICLR)}.

\bibitem[Wang et al.(2022)]{wang2022interpretability}
Wang, K., Variengien, A., Conmy, A., Shlegeris, B., \& Steinhardt, J.
(2022).
Interpretability in the Wild: a Circuit for Indirect Object Identification in GPT-2
Small.
\emph{arXiv preprint arXiv:2211.00593}.

\end{thebibliography}

\end{document}
